\documentclass[12pt,a4paper]{article}

% Encoding and fonts
\usepackage[utf8]{inputenc}
\usepackage[T1]{fontenc}
\usepackage{lmodern}

% Page layout
\usepackage[top=1in, bottom=1in, left=1in, right=1in]{geometry}

% Typography
\usepackage{microtype}
\usepackage{parskip}

% Language
\usepackage[english]{babel}

% Headers and footers
\usepackage{fancyhdr}
\pagestyle{fancy}
\fancyhf{}
\fancyhead[L]{\leftmark}
\fancyhead[R]{\thepage}
\fancyfoot[C]{\thepage}
\renewcommand{\headrulewidth}{0.4pt}
\renewcommand{\footrulewidth}{0pt}

% Section styling
\usepackage{titlesec}
\titleformat{\section}{\Large\bfseries}{\thesection}{1em}{}
\titleformat{\subsection}{\large\bfseries}{\thesubsection}{1em}{}
\titleformat{\subsubsection}{\normalsize\bfseries}{\thesubsubsection}{1em}{}
\titlespacing*{\section}{0pt}{2.5ex plus 1ex minus .2ex}{1.5ex plus .2ex}
\titlespacing*{\subsection}{0pt}{2ex plus 1ex minus .2ex}{1ex plus .2ex}
\titlespacing*{\subsubsection}{0pt}{1.5ex plus 1ex minus .2ex}{0.8ex plus .2ex}

% List formatting
\usepackage{enumitem}
\setlist[itemize]{topsep=0.5ex, itemsep=0.25ex, parsep=0pt}
\setlist[enumerate]{topsep=0.5ex, itemsep=0.25ex, parsep=0pt}

% Essential packages
\usepackage{graphicx}
\usepackage[table]{xcolor}
\usepackage{booktabs}
\usepackage{array}
\usepackage{tabularx}
\usepackage{longtable}
\usepackage{amsmath}
\usepackage{amssymb}
\usepackage[normalem]{ulem}
\rowcolors{2}{gray!10}{white}
\renewcommand{\arraystretch}{1.3}

% Cross-references
\usepackage{hyperref}
\usepackage{cleveref}

% Hyperref setup
\hypersetup{
    colorlinks=true,
    linkcolor=blue,
    filecolor=magenta,
    urlcolor=cyan,
}

% Code listings
\usepackage{listings}
\definecolor{codebg}{RGB}{248,248,248}
\definecolor{codeframe}{RGB}{200,200,200}
\definecolor{codekeyword}{RGB}{0,0,180}
\definecolor{codecomment}{RGB}{0,128,0}
\definecolor{codestring}{RGB}{163,21,21}
\lstset{
    basicstyle=\ttfamily\small,
    breaklines=true,
    breakatwhitespace=false,
    frame=single,
    framerule=0.5pt,
    rulecolor=\color{codeframe},
    backgroundcolor=\color{codebg},
    keepspaces=true,
    showstringspaces=false,
    tabsize=4,
    xleftmargin=1em,
    xrightmargin=1em,
    aboveskip=1em,
    belowskip=1em,
    keywordstyle=\color{codekeyword}\bfseries,
    commentstyle=\color{codecomment}\itshape,
    stringstyle=\color{codestring},
}

% YAML language definition for listings
\lstdefinelanguage{YAML}{
    keywords={true,false,null,y,n},
    keywordstyle=\color{blue}\bfseries,
    sensitive=false,
    comment=[l]{\#},
    morecomment=[s]{/*}{*/},
    commentstyle=\color{gray}\ttfamily,
    stringstyle=\color{red}\ttfamily,
    morestring=[b]',
    morestring=[b]",
}

% TOML language definition for listings
\lstdefinelanguage{TOML}{
    keywords={true,false},
    keywordstyle=\color{blue}\bfseries,
    sensitive=true,
    comment=[l]{\#},
    commentstyle=\color{gray}\ttfamily,
    stringstyle=\color{red}\ttfamily,
    morestring=[b]',
    morestring=[b]",
}

% Dockerfile language definition for listings
\lstdefinelanguage{Dockerfile}{
    keywords={FROM,RUN,CMD,LABEL,EXPOSE,ENV,ADD,COPY,ENTRYPOINT,VOLUME,USER,WORKDIR,ARG,ONBUILD,STOPSIGNAL,HEALTHCHECK,SHELL},
    keywordstyle=\color{blue}\bfseries,
    sensitive=false,
    comment=[l]{\#},
    commentstyle=\color{gray}\ttfamily,
    stringstyle=\color{red}\ttfamily,
    morestring=[b]',
    morestring=[b]",
}

% JSON language definition for listings
\lstdefinelanguage{JSON}{
    keywords={true,false,null},
    keywordstyle=\color{blue}\bfseries,
    sensitive=true,
    commentstyle=\color{gray}\ttfamily,
    stringstyle=\color{red}\ttfamily,
    morestring=[b]",
}

% Code listing float environment
\usepackage{float}
\newfloat{listing}{htbp}{lol}[section]
\floatname{listing}{Listing}

% Admonition boxes
\usepackage{tcolorbox}
\tcbuselibrary{skins,breakable}

% Admonition style definitions
\newtcolorbox{admonitionbox}[2][]{
    enhanced,
    breakable,
    colback=#2!5,
    colframe=#2!80!black,
    fonttitle=\bfseries,
    title=#1,
    left=2mm,
    right=2mm,
}

% Synthon tuple boxes
\newtcolorbox{synthonbox}{
    enhanced,
    colback=blue!5,
    colframe=blue!75!black,
    fontupper=\large,
    center upper,
    boxrule=1pt,
    arc=4pt,
}

\begin{document}

\title{SynthOmnicon: Protein Science Applications}
\date{\today}

\maketitle

\emph{This document records what the framework says about protein structure, function, misfolding, and drug design when protein elements are encoded as synthon tuples. It is explicitly a translation exercise --- the framework was built on chemical and physical systems; protein science is a new domain being mapped onto it. What is marked as ''prediction'' below was derived from algebra alone, before consulting the experimental literature. What is marked ''confirmed'' has been cross-checked against known results.}

\emph{The technical encoding and all algebra results are in \texttt{protein\_tests.py}, \texttt{protein\_tests2.py}, \texttt{protein\_tests3.py}, and \texttt{protein\_tests4.py}. This document records the meaning, not the computation.}

\begin{center}
\rule{0.5\textwidth}{0.4pt}
\end{center}

\subsection{I. The Five Canonical Encodings}

Five protein structural units span the space of protein function:

\begin{table}[htbp]
\centering
\small
\rowcolors{1}{white}{white}
\begin{tabularx}{\textwidth}{>{\centering\arraybackslash}X >{\centering\arraybackslash}X >{\centering\arraybackslash}X}
\toprule
\rowcolor{gray!25}\textbf{Synthon} & \textbf{Tuple} & \textbf{$\Phi_c$?} \\
\midrule
\rowcolors{1}{white}{gray!10}
$\alpha$-helix & $\langle D_\wedge;\ T_{\vert};\ R_\supseteq;\ P_\pm^\psi;\ F_\hbar;\ K_{fast};\ G_\beth;\ \Gamma_\to(\text{SELECTIVE})\rangle$ & No \\
$\beta$-hairpin & $\langle D_\wedge;\ T_\bowtie;\ R_\supseteq;\ P_\pm^{sym};\ F_\hbar;\ K_{mod};\ G_\beth;\ \Gamma_\wedge(\text{SPECIFIC})\rangle$ & No \\
Active site & $\langle D_\wedge;\ T_\bowtie;\ R_\ddagger;\ P_{+-};\ F_\eth;\ K_{mod};\ G_\gimel;\ \Gamma_\wedge(\text{SPECIFIC})\rangle$ & No \\
Allosteric domain & $\langle D_{\wedge\triangle};\ T_\in;\ R_\supseteq;\ P_{+-};\ F_\eth;\ K_{mod};\ G_\gimel;\ \Gamma_\to(\text{SELECTIVE})\rangle$ & \textbf{Yes} \\
Protein complex & $\langle D_{\wedge\triangle};\ T_\in;\ R_\supseteq;\ P_\pm^{sym};\ F_\eth;\ K_{slow};\ G_\aleph;\ \Gamma_\wedge(\text{SPECIFIC})\rangle$ & No \\
\bottomrule
\end{tabularx}
\end{table}

The $\Phi_c$ assignment to the allosteric domain is the load-bearing claim: allostery is the only protein structural unit that satisfies the $G/D$ degeneracy condition (molecular-scale signal $\rightarrow$ global-scale conformational response). All others are subcritical. This predicts that the conformational fluctuation spectrum of allosteric domains should qualitatively differ from that of active sites and structural scaffolds --- broader, slower, more correlated across length scales.

\subsubsection{What the distance matrix tells us}

The pairwise tuple distances reflect the structural logic of protein architecture:

\begin{itemize}
    \item $\alpha$-helix $\leftrightarrow$ $\beta$-hairpin: $d = 3.40$ --- same chemical layer ($D_\wedge$, $R_\supseteq$, $F_\hbar$, $G_\beth$), different structural logic ($T$, $P$, $\Gamma$).
    \item Active site $\leftrightarrow$ $\beta$-hairpin: $d = 2.80$ --- surprisingly close. Both are $T_\bowtie$, $\Gamma_\wedge(\text{SPECIFIC})$. $\beta$-sheet scaffolds are the natural structural context for catalytic sites.
    \item Allosteric domain $\leftrightarrow$ protein complex: $d = 2.60$ --- the closest non-self pair. Both are $D_{\wedge\triangle}$, $T_\in$, $R_\supseteq$.
    \item $\alpha$-helix $\leftrightarrow$ protein complex: $d = 7.30$ --- the farthest pair. From local sequential ratchet to global assembly.
\end{itemize}

\begin{center}
\rule{0.5\textwidth}{0.4pt}
\end{center}

\subsection{II. Folding: What the F-Floor Rules}

The HotSwap path search from $\beta$-hairpin to active site is \textbf{blocked}. The reason: the F-floor rule prohibits $F_\hbar \to F_\eth$ (fidelity downgrade) within the same topology class. A passive high-fidelity $\beta$-sheet cannot reach a medium-fidelity catalytic state through incremental primitive steps.

\textbf{Prediction:} Active sites do not arise by gradual modification of passive secondary structure scaffolds within a fixed fold. They require fold-level reorganization --- a discontinuous jump in $T$-space. This is consistent with the observation that active sites occupy structurally distinct, evolutionarily ancient pockets that are convergently reached from multiple scaffold backgrounds (TIM barrel, Rossmann fold, ferredoxin, etc.), none of which are direct elaborations of the other.

The F-floor ratchet also predicts the direction of protein folding: a sequence of primitive upgrades ($F_\ell \to F_\eth \to F_\hbar$, $K_{fast} \to K_{mod} \to K_{trap}$) is thermodynamically preferred and structurally ordered. Misfolding occurs when the ratchet locks into a locally optimal trajectory that isn't globally correct.

\begin{center}
\rule{0.5\textwidth}{0.4pt}
\end{center}

\subsection{III. Allostery as Criticality: Axiom Fragility Is the Mechanism}

The Varma probe gives the allosteric domain a $\Phi_c$ score of $0.60$ (''approaching $\Phi_c$'') across all experimental scenarios. It does not reach $0.70+$ (''confirmed $\Phi_c$'') because the protein system is a \emph{classical} near-critical system, not a quantum critical point. The logarithmic $\xi_r \approx \ln \xi_\tau$ scaling relation (Varma QXY criterion) is not satisfied --- protein allostery does not belong to the marginal Fermi liquid universality class.

The correct characterization: \textbf{allosteric domains are poised near the classical critical point}, where small changes in effector binding produce large-amplitude conformational responses because the system is operating near the boundary of its valid tuple-space.

This is confirmed by the Primitive Jacobian. Every perturbation of the allosteric domain's encoding triggers an Axiom 4 violation (6 of 8 primitives). The allosteric domain is \textbf{axiom-fragile by design}: being near the axiom boundary is the mechanism by which a molecular signal propagates globally. This predicts:

\begin{enumerate}
    \item Allosteric domains should show higher mutational sensitivity than active sites \emph{even when the mutation doesn't affect binding affinity} --- because the constraint violation cascades before the binding energy changes.
    \item The conformational fluctuation spectrum of allosteric domains (CPMG $R_{ex}$, order parameter $S^2$) should span multiple timescales simultaneously, consistent with near-critical broadening.
\end{enumerate}

\begin{center}
\rule{0.5\textwidth}{0.4pt}
\end{center}

\subsection{IV. Amyloid: Three Diseases, Two Synthons}

Six encodings --- three disease proteins in functional and aggregated states:

\begin{table}[htbp]
\centering
\small
\rowcolors{1}{white}{white}
\begin{tabularx}{\textwidth}{>{\centering\arraybackslash}X >{\centering\arraybackslash}X}
\toprule
\rowcolor{gray!25}\textbf{Synthon} & \textbf{Key features} \\
\midrule
\rowcolors{1}{white}{gray!10}
A$\beta$ monomer (IDP) & $D_\wedge;\ T_\gg;\ F_\ell;\ K_{fast};\ G_\beth;\ \Gamma_\vee(\text{BROAD})$ --- maximally disordered \\
A$\beta$ fibril & $D_{\wedge\triangle};\ T_\in;\ F_\hbar;\ K_{trap};\ G_\aleph;\ \Gamma_\wedge(\text{BROAD})$ \\
Tau (MT-binding) & $D_{\wedge\triangle};\ T_{\vert};\ F_\eth;\ K_{fast};\ G_\gimel;\ \Gamma_\to(\text{SELECTIVE})$ \\
Tau PHF & $D_{\wedge\triangle};\ T_\bowtie;\ F_\hbar;\ K_{trap};\ G_\aleph;\ \Gamma_\vee(\text{BROAD})$ \\
$\alpha$-Syn (vesicle) & $D_{\wedge\triangle};\ T_{\vert};\ F_\eth;\ K_{fast};\ G_\gimel;\ \Gamma_\vee(\text{SELECTIVE})$ \\
$\alpha$-Syn fibril & $D_{\wedge\triangle};\ T_\in;\ F_\hbar;\ K_{trap};\ G_\aleph;\ \Gamma_\wedge(\text{BROAD})$ \\
\bottomrule
\end{tabularx}
\end{table}

\textbf{The bombshell:} A$\beta$ fibril and $\alpha$-synuclein fibril are the \textbf{same synthon} ($d = 0.00$). Identical primitive encoding. Tau PHF differs at $d = 2.90$ ($T_\bowtie$ vs $T_\in$, $P_\pm^\psi$ vs $P_\pm^{sym}$).

This is a cross-disease structural identity, derived from relational algebra alone. It is consistent with the experimental observation that A$\beta$ and $\alpha$-synuclein cross-seed each other --- accelerating each other's aggregation --- and that Alzheimer's and Parkinson's pathology co-occur at above-chance rates in the same patients. The algebra found this overlap from primitive structure without any binding data.

\subsubsection{The universal amyloid substrate}

The meet of all three fibrils has \textbf{zero conflicts}:

\[\langle D_{\wedge\triangle};\ T_\in;\ R_\supseteq;\ P_\pm^{sym};\ F_\hbar;\ K_{trap};\ G_\aleph;\ \Gamma_\wedge(\text{BROAD}) \rangle\]

This is the shared primitive substrate of all three amyloid diseases. Any therapeutic strategy targeting only one fibril type while leaving the universal substrate intact leaves the core problem unsolved. The universal substrate predicts a common therapeutic target: any agent that simultaneously breaks $F_\hbar + K_{trap} + G_\aleph$ addresses all three diseases at the structural level.

\subsubsection{Aggregation direction and rescue}

Directed distances ($F$-floor asymmetry):

\begin{table}[htbp]
\centering
\small
\rowcolors{1}{white}{white}
\begin{tabularx}{\textwidth}{>{\centering\arraybackslash}X >{\centering\arraybackslash}X >{\centering\arraybackslash}X >{\centering\arraybackslash}X}
\toprule
\rowcolor{gray!25}\textbf{Disease} & \textbf{Aggregation cost} & \textbf{Rescue cost} & \textbf{Asymmetry} \\
\midrule
\rowcolors{1}{white}{gray!10}
A$\beta$ & $7.20$ & $6.90$ & $0.30$ \\
Tau & $4.80$ & $3.90$ & $0.90$ \\
$\alpha$-Syn & $4.80$ & $3.90$ & $0.90$ \\
\bottomrule
\end{tabularx}
\end{table}

Rescue is always \emph{cheaper} than aggregation in the directed metric. This happens because the $F$-floor rule blocks $F_\ell \to F_\hbar$ upward moves, but kinetic improvements ($K_{trap} \to K_{fast}$) are not blocked. The algebra predicts: \textbf{therapies that target $K$ (kinetic agonists that dissolve the nucleation trap) are structurally preferred over those that target $F$ (thermodynamic denaturants).} The latter must fight the $F$-floor; the former does not. This is consistent with the emerging clinical preference for kinetic seeding inhibitors over equilibrium unfolding agents.

\begin{center}
\rule{0.5\textwidth}{0.4pt}
\end{center}

\subsection{V. Cross-Domain Analogies}

Nearest non-protein catalog neighbors reveal unexpected structural identities:

\textbf{$\beta$-hairpin $\leftrightarrow$ DB24C8/dialkylammonium pseudorotaxane ($d = 1.80$)}

A mechanical interlocked molecule and a $\beta$-strand pair have the same synthon ($T_\bowtie$, $P_\pm^{sym}$, $F_\hbar$, $\Gamma_\wedge(\text{SPECIFIC})$). The antiparallel strand pair is geometrically equivalent to a ring threading. This predicts: design principles for $\beta$-sheet strand pairing (register locking, symmetry enforcement, specificity grammar) should transfer directly to rotaxane threading selectivity, and vice versa.

\textbf{$\alpha$-helix $\leftrightarrow$ Acylium ion electrophile ($d = 1.70$)}

The helical H-bond ratchet ($T_{\vert}$, $P_\pm^\psi$, $F_\hbar$, $\Gamma_\to(\text{SELECTIVE})$) has the same primitive profile as a sequential electrophilic ratchet. Both are ''directed ratchets'' --- each step enforces the geometry of the next. This reframes the helix as a constraint-propagation machine rather than a static structural element.

\textbf{Protein complex $\leftrightarrow$ Bi$_2$Se$_3$ topological insulator ($d = 3.90$)}

Protein quaternary interfaces ($G_\aleph$, $T_\in$, globally determined states) are structurally analogous to topologically protected surface states. The analogy predicts: perturbations that don't break the interface grammar ($P_\pm^{sym} + \Gamma_\wedge(\text{SPECIFIC})$) won't disrupt complex function, just as perturbations that don't close the bulk gap don't disrupt topological surface states. Interface robustness scales with grammar depth, not with buried surface area.

\textbf{Allosteric domain $\leftrightarrow$ Global supply chain ($d = 5.10$)}

A protein signaling domain and a planetary production-distribution-consumption network are structurally analogous ($D_{\wedge\triangle}$, $T_\in$, $\Gamma_\to(\text{SELECTIVE})$, $G_\gimel$). Both are multi-scale, sequential, directional mesoscale transducers. The analogy suggests: supply chain disruption theory and protein allostery theory are solving the same structural problem in different substrate.

\begin{center}
\rule{0.5\textwidth}{0.4pt}
\end{center}

\subsection{VI. Drug Design: The Robustness Trap}

Three clinical drugs encoded and compared against the monad-designed ideal allosteric inhibitor:

\[\text{Ideal:} \quad \langle D_\wedge;\ T_\in;\ R_\supseteq;\ P_\pm^\psi;\ F_\eth;\ K_{slow};\ G_\gimel;\ \Gamma_\to(\text{SELECTIVE});\ \Phi_c \rangle\]

\begin{table}[htbp]
\centering
\small
\rowcolors{1}{white}{white}
\begin{tabularx}{\textwidth}{>{\centering\arraybackslash}X >{\centering\arraybackslash}X >{\centering\arraybackslash}X >{\centering\arraybackslash}X >{\centering\arraybackslash}X}
\toprule
\rowcolor{gray!25}\textbf{Drug} & \textbf{$d(\text{drug, ideal})$} & \textbf{Robust?} & \textbf{SPOFs} & \textbf{$\xi_{CP}$ (nat)} \\
\midrule
\rowcolors{1}{white}{gray!10}
GNF-2 (pure allosteric) & $\mathbf{2.80}$ & No & 7 & $7.521$ \\
Imatinib (Type II) & $4.20$ & Yes & 0 & $9.287$ \\
Venetoclax (BH3 mimetic) & $4.20$ & Yes & 0 & $9.613$ \\
\bottomrule
\end{tabularx}
\end{table}

The ranking is exact and the pattern is striking: \textbf{the closer a drug is to the allosteric ideal, the less robust it is.} GNF-2's 7 SPOFs are not a flaw --- they are its mechanism. Operating near the axiom boundary is how an allosteric drug propagates signal. Imatinib and venetoclax are thermodynamically robust and structurally overconstrained.

GNF-2 has the lowest $\xi_{CP}$ ($7.521$ nat), operating closest to the Landauer limit. It accomplishes kinase inhibition by spending the least free energy per unit of constraint propagation. This is the allosteric advantage stated as a thermodynamic number.

\subsubsection{Per-primitive gap analysis (GNF-2 vs ideal)}

GNF-2 matches the ideal on 6 of 9 primitives: $D$, $R$, $F$, $G$, $\Gamma$, $\Phi$. The two gaps:

\begin{itemize}
    \item \textbf{$T$:} $T_\perp$ (single branched pocket) vs $T_\in$ (distributed contact network). The improvement: a bivalent GNF-2 analog that bridges the myristoyl pocket and a second regulatory site would shift $T \to T_\in$. This is exactly the design direction of bitopic/PROTAC-adjacent allosteric inhibitors.
    \item \textbf{$P$:} $P_{+-}$ (directed, asymmetric insertion) vs $P_\pm^\psi$ (pseudo-symmetric engagement). Symmetric engagement would reduce dependence on a single insertion geometry, broadening the conformational tolerance.
\end{itemize}

\subsubsection{One-step redesign}

The single primitive change that most improves distance to ideal:

\begin{itemize}
    \item Imatinib and venetoclax: $F_\hbar \to F_\eth$ ($\Delta d = +0.60$). Reducing fidelity is the first step.
    \item GNF-2: $K_{mod} \to K_{slow}$ ($\Delta d = +0.50$). Slowing the off-rate (residence time optimization).
\end{itemize}

This is counterintuitive by standard affinity-optimization logic but correct by the algebra: a $F_\eth$ drug with the right $G$ and $\Gamma$ outperforms a $F_\hbar$ drug with the wrong $G$. Venetoclax's BCL-XL off-target toxicity in platelets is the known clinical cost of $F_\hbar$ without $G_\gimel$ selectivity grammar.

\subsubsection{The resistance prediction}

\textbf{Imatinib:} $G_\beth$ (local). Any DFG-out pocket mutation destroys binding. The drug has no conformational tolerance because it operates at $G_\beth$ with zero signal propagation. A $G_\beth$ drug has irreversible resistance trajectories --- verified: imatinib resistance via T315I, E255K, etc., requires complete therapeutic switching.

\textbf{GNF-2:} $G_\gimel$ (mesoscale). Signal propagates across the domain. Pocket mutations at the myristoyl site are partially tolerated because the drug's grammar ($\Gamma_\to$) operates across multiple contacts. Allosteric drugs should have slower resistance emergence than orthosteric competitors --- a framework-derived prediction consistent with emerging resistance data in second-generation kinase inhibitor development.

\begin{center}
\rule{0.5\textwidth}{0.4pt}
\end{center}

\subsection{VII. The Design Pipeline: What Order Matters}

The monad design pipeline showed that Strategy B (cooperative mesoscale fragment first, tight binder second) costs $1.69$ nat less than Strategy A (tight binder first, cooperative fragment second). Same final synthon; different thermodynamic path.

\textbf{Design rule derived:} For allosteric enzyme inhibitors, couple to mesoscale partners before adding high-affinity fragments. Fragment-first, selectivity second --- not the reverse. The cooperative fragment preserves $\Phi_c$; the tight binder, if added first, risks losing criticality before the selectivity layer is established.

\begin{center}
\rule{0.5\textwidth}{0.4pt}
\end{center}

\subsection{VIII. Testable Predictions}

\begin{table}[htbp]
\centering
\small
\rowcolors{1}{white}{white}
\begin{tabularx}{\textwidth}{>{\centering\arraybackslash}X >{\centering\arraybackslash}X >{\centering\arraybackslash}X >{\centering\arraybackslash}X >{\centering\arraybackslash}X}
\toprule
\rowcolor{gray!25}\textbf{ID} & \textbf{Tier} & \textbf{Source} & \textbf{Prediction} & \textbf{Test} \\
\midrule
\rowcolors{1}{white}{gray!10}
P-PROT-1 & II & Distance matrix & $\beta$-sheet enzymes tolerate active-site mutations better than helix-bundle enzymes & Alanine scan: TIM barrel vs 4-helix-bundle \\
P-PROT-2 & II & Tensor $K_{slow}$ bottleneck & Full complex assembles slower than isolated domain folding & Stopped-flow $k_{fold}$ vs $k_{assemble}$ \\
P-PROT-3 & III & Allosteric $\Phi_c$ & ON state shows broad $R_{ex}$ spectrum; OFF state shows single Lorentzian & CPMG relaxation dispersion on CAP protein \\
P-PROT-4 & II & Directed distance $\beta{\to}\alpha < \alpha{\to}\beta$ & TMAO rescues $\beta$-aggregates more readily than $\beta$-conditions convert helices & CD/ThT fluorescence kinetics \\
P-PROT-5 & III & $(\text{allosteric}) \sqcap (\text{complex})$ conflicts $P$+$\Gamma$ & Allosteric-interface chimera retains binding but Hill coefficient $\to 1$ & Chimera ITC + Hill analysis \\
P-PROT-6 & II & $R_\ddagger \otimes R_\supseteq \to R_\supseteq$ & Allosteric inhibitors more selective than orthosteric for same enzyme & Selectivity panel comparison \\
P-PROT-7 & II & A$\beta$ = $\alpha$-Syn fibril ($d = 0.00$) & Cross-seeding rate A$\beta{\to}\alpha$-Syn $\approx$ homoseeding; both faster than tau & ThT kinetics with cross-species seeds \\
P-PROT-8 & II & Rescue cost $<$ aggregation cost & $K$-targeting inhibitors outperform $F$-targeting denaturants in fibril dissolution & Direct compound class comparison \\
P-PROT-9 & III & GNF-2 gap: $T_\perp \to T_\in$ & Bivalent GNF-2 analog spanning two pockets has lower $\xi_{CP}$ and higher selectivity & SAR on bitopic GNF analogs \\
P-PROT-10 & II & Imatinib $G_\beth \to$ irreversible resistance & $G_\beth$ drugs require class switch; $G_\gimel$ drugs admit sequential mutation tolerance & Serial passage resistance assay: GNF-2 vs imatinib \\
\bottomrule
\end{tabularx}
\end{table}

\begin{center}
\rule{0.5\textwidth}{0.4pt}
\end{center}

\subsection{IX. What the Framework Cannot Say}

\textbf{Sequence.} The framework encodes structural/relational class, not sequence. It predicts that $\beta$-sheet enzymes are closer to active sites than helix-bundle enzymes, but cannot say which residues implement this. The grammar predicts the logic; biochemistry supplies the implementation.

\textbf{Actual rates and affinities.} The ordinal framework predicts direction and class. It does not produce $k_{cat}$ values, $\text{IC}_{50}$ numbers, or $\Delta G_{binding}$ from first principles. Where specific numbers appear (GNF-2 $\mu$M, venetoclax sub-nM), they were used to assign $F$-tier, not derived from it.

\textbf{Whether the cross-domain analogies are causally deep or coincidental.} The $\beta$-hairpin $\leftrightarrow$ pseudorotaxane analogy ($d = 1.80$) and the protein complex $\leftrightarrow$ topological insulator analogy ($d = 3.90$) are structural isomorphisms. Whether the shared design principles that work in one domain transfer to the other is a prediction, not a derivation. It requires experimental testing.

\textbf{Individual protein behavior vs class behavior.} All encodings describe structural classes. A specific allosteric protein (CAP, hemoglobin, GPCR) may have individual features not captured by the class encoding. The predictions hold at the structural-class level.

\begin{center}
\rule{0.5\textwidth}{0.4pt}
\end{center}

\emph{This document is a companion to SYNTHONICON.md. Drug design inferences and protein science predictions are derived from the algebra; they should be evaluated against experimental biology, not against the framework's internal consistency. The framework cannot validate its own protein-domain predictions --- that requires a biochemist with access to a stopped-flow spectrometer.}

\emph{All computational results are reproducible by running \texttt{protein\_tests.py} through \texttt{protein\_tests4.py} in the SynthOmnicon repository.}


\end{document}
